\section{Project Statement}

\subsection{Project Title}

\textbf{Multi-Agent AI System for Structured Data Collaboration and Natural Language Interaction}

\subsection{Statement of Intent}

This project proposes the design and development of \projectname{}, a lightweight but credible multi-agent AI platform that helps people and software systems collaborate around structured data. The system accepts JSON, YAML, and CSV inputs, interprets natural language instructions, coordinates specialized agents, invokes validated data tools through the Model Context Protocol (MCP), retrieves trusted schema and transformation knowledge through retrieval-augmented generation (RAG), and returns correct outputs with clear explanations.

The goal is not to build another file converter. The goal is to build an intelligent structured-data operations layer: a system that can convert, validate, explain, audit, and ask for human approval when a transformation could change meaning. This direction fits the HuReDee-IoYou AI OJT program because it demonstrates practical backend engineering, model integration, agent orchestration, protocol design, deployment readiness, and product thinking in one cohesive prototype.

For the medical and healthcare application track, \projectname{} should be treated as clinical data support infrastructure, not an autonomous diagnosis system. It should help transform, validate, retrieve, summarize, and explain healthcare data while preserving provenance, uncertainty, limitations, human accountability, and clinician review.

\subsection{Core Objectives}

\begin{itemize}
  \item Design a modular multi-agent architecture for structured-data workflows.
  \item Support reliable conversion among JSON, YAML, and CSV while preserving schema intent.
  \item Enable natural language control for conversion, validation, cleaning, and explanation tasks.
  \item Use RAG to ground decisions in approved schemas, transformation examples, and data dictionaries.
  \item Apply advanced retrieval and self-supervised representation learning to improve schema matching, medical terminology alignment, and evidence retrieval.
  \item Expose deterministic data tools through MCP servers so agents use standard, auditable interfaces.
  \item Add clinician/human-in-the-loop checkpoints for ambiguous, destructive, sensitive, or medically consequential operations.
  \item Provide healthcare-grade explainability: evidence traces, uncertainty, source provenance, data-quality warnings, and interpretable user-facing reports.
  \item Deliver a deployable API-driven prototype with documentation, tests, evaluation metrics, and a demonstration scenario.
\end{itemize}

\subsection{Team Vision}

The team vision is to bridge two worlds that usually stay separated: strict machine-readable data and natural human collaboration. \projectname{} should let a non-technical user ask, ``Clean this CSV and explain the anomalies,'' while the system behaves like disciplined software: parsing exactly, validating against explicit rules, preserving provenance, and refusing risky actions until a person approves them.
